\chapter{Background}
\label{chap:background}

\section{Literature Review}
\label{sec:literature_review}

\subsection{Thesis Topic}
\label{sec:thesis_topic}
RGB Images Generation Based on Feature Detection Techniques: The Case of SegNet

\subsection{Introduction}
\label{sec:lit_intro}
Brain tumor segmentation in Magnetic Resonance Imaging (MRI) has long been recognized as a critical component in the diagnosis, treatment planning, and outcome evaluation for patients with brain tumors. Manual delineation of tumor regions by radiologists, while historically the gold standard, is labor intensive, subject to inter-observer variability, and often impractical for high-throughput clinical environments. With the increasing availability of large annotated datasets such as those provided by the Brain Tumor Segmentation (BraTS) challenges \cite{menze2015brats,bakas2017brats,brats2020}, there has been growing interest in automating this process through advanced computational methods.

Recent years have witnessed a paradigm shift from classical image processing techniques towards data-driven approaches, particularly deep learning methodologies. Convolutional Neural Networks (CNNs) have emerged as the most promising tool due to their capacity to learn complex, hierarchical feature representations directly from data. However, while several CNN-based architectures have achieved impressive accuracy—most notably U-Net \cite{ronneberger2015unet}—challenges still persist in balancing segmentation accuracy with computational and memory efficiency.

This literature review traces the evolution of techniques for brain tumor MRI segmentation. It begins with early classical methods, covers the transition into learning-based techniques, and progressively narrows the focus towards deep learning encoder–decoder architectures. The review concludes with a detailed examination of the SegNet architecture \cite{badrinarayanan2017segnet}, its advantages and limitations, and the modifications applied by recent studies to enhance its performance. This survey establishes a foundation for the proposed thesis approach, which integrates feature detection techniques with a SegNet-based framework to generate RGB segmentation maps.

\subsection{Traditional Methods in Brain Tumor Segmentation}
\label{sec:traditional_methods}

\subsubsection{Classical Image Processing Approaches}
\label{sec:classical_image_processing}
Prior to the advent of deep learning, brain tumor segmentation largely depended on classical image processing methods. Techniques such as thresholding, region growing, and edge detection were extensively applied to MRI scans. Threshold-based methods leveraged differences in image intensity between tumor and normal tissue, but often fell short in cases of overlapping intensity ranges and image noise. Region growing algorithms attempted to expand from a user-defined seed point based on pixel similarity, but they struggled with heterogeneous tumor structures and irregular boundaries. The application of edge detection, typically using Sobel \cite{sobel1968isotropic} or Canny operators, helped in identifying high-gradient regions corresponding to tissue boundaries, yet these methods were sensitive to noise and often resulted in fragmented segmentation outputs \cite{gordillo2013state}.

\subsubsection{Machine Learning with Handcrafted Features}
\label{sec:handcrafted_ml}
To overcome the limitations of purely heuristic methods, researchers began to apply traditional machine learning techniques to MRI segmentation. In such approaches, specific features—such as texture, intensity gradients, and local binary patterns (LBP)—were manually engineered from the MRI data. These features were then classified using algorithms like support vector machines (SVMs), random forests, or k-means clustering. Although these methods offered improved segmentation quality in some instances, they were inherently limited by the quality of the handcrafted features and often required extensive post-processing (e.g., Conditional Random Fields, CRFs) to improve spatial consistency. Despite certain successes, the performance of these methods was unsatisfactory for clinical deployment due to issues with reproducibility and robustness across different imaging protocols.

\subsection{Transition to Deep Learning-Based Segmentation}
\label{sec:transition_dl}

\subsubsection{Early Applications of CNNs}
\label{sec:early_cnns}
With the advent of CNNs, segmentation methods saw a notable improvement. Early CNNs were trained on small image patches, allowing the network to learn local features—such as edges and textures—that were later aggregated to determine if a region belonged to a tumor. This patch-based approach marked a significant departure from manual feature engineering. Studies demonstrated that CNNs could outperform classical methods by learning discriminative representations directly from raw data, even when trained on relatively small datasets \cite{litjens2017survey}.

\textbf{Analogy:} Imagine a CNN as a team of specialists where each member learns to detect specific patterns (e.g., edges, textures). Together, they compile these observations to create a comprehensive map of the tumor region, much like how a committee of experts might combine their assessments to make a conclusive diagnosis.

\subsubsection{Fully Convolutional Networks (FCNs)}
\label{sec:fcns}
A major breakthrough in segmentation was achieved with Fully Convolutional Networks (FCNs), which replaced fully connected layers with convolutional layers to output pixel-wise predictions for entire images. This end-to-end training enabled models to process full images in one pass rather than in a patch-wise manner, significantly accelerating inference. However, early FCN models suffered from coarse segmentation outputs due to extensive downsampling by pooling layers, which reduced spatial resolution and blurred tumor boundaries.

\subsubsection{Emergence of Encoder–Decoder Architectures}
\label{sec:encoder_decoder_emergence}
To address the resolution loss in FCNs, encoder–decoder architectures were developed. In these models, the encoder extracts hierarchical features while progressively reducing spatial dimensions, and the decoder upscales these low-resolution features back to the original image size to generate detailed segmentation maps. U-Net, introduced by Ronneberger, Fischer, and Brox \cite{ronneberger2015unet}, became a seminal work in this area. U-Net is distinguished by its symmetric design and the use of skip connections that transfer high-resolution features from the encoder to the decoder, significantly aiding in the accurate delineation of tumor boundaries.

Despite U-Net's high accuracy, its reliance on storing full-resolution feature maps can be resource-intensive, leading researchers to explore alternatives that reduce memory consumption without compromising performance.

\subsection{Encoder–Decoder Architectures: U-Net vs. SegNet}
\label{sec:unet_vs_segnet}

\subsubsection{U-Net and Its Dominance in Medical Imaging}
\label{sec:unet_dominance}
U-Net has emerged as the gold standard for various biomedical segmentation tasks. Its design allows for the combination of contextual information and precise localization. Specifically, U-Net's skip connections recycle detailed feature maps from the encoder, enabling the decoder to produce finely detailed segmentation masks. Studies applying U-Net to the BraTS datasets have reported high Dice similarity coefficients, often exceeding 0.85 for whole tumor segmentation \cite{ronneberger2015unet}. However, U-Net's requirement for extensive memory to store these skip connections can be a considerable drawback in real-time applications and on systems with limited GPU resources.

\subsubsection{SegNet: Design, Advantages, and Limitations}
\label{sec:segnet_design}

\paragraph{Architectural Overview}
\label{sec:segnet_overview}
SegNet, as introduced by Badrinarayanan, Kendall, and Cipolla \cite{badrinarayanan2017segnet}, shares the encoder–decoder framework but introduces a key innovation: the use of pooling index transfer. Its encoder is identical to the convolutional layers of the VGG16 network, where each max-pooling operation not only reduces spatial dimensions but also records the indices of maximal activations. In the decoder, these stored indices guide the unpooling process, efficiently restoring spatial resolution.

\begin{figure}[ht]
    \centering
    % Insert a diagram of the SegNet architecture here
    \caption{A schematic diagram of SegNet illustrating the encoder's max pooling layers, the storage of pooling indices, and the decoder's unpooling process using these indices.}
    \label{fig:segnet_architecture}
\end{figure}

\paragraph{Pooling Index Transfer Mechanism}
\label{sec:pooling_indices}
The pooling index transfer mechanism lies at the heart of SegNet's efficiency. During encoding, the max-pooling operation condenses feature maps while retaining the spatial locations of the most significant activations. In the decoder, these indices are used to "unpool" the feature maps, essentially placing the activations back into their original spatial locations. This step is followed by standard convolution to fill in gaps, thus producing a dense segmentation map without the computational overhead of storing entire feature maps from the encoder.

\textbf{Analogy:} Think of the encoder as leaving behind "breadcrumbs" (pooling indices) that mark where important features were found. The decoder uses these breadcrumbs to reconstruct the image accurately, ensuring that critical boundary information is not lost in the pooling process.

\paragraph{Comparative Advantages}
\label{sec:segnet_advantages}
SegNet offers several distinct advantages:
\begin{itemize}
    \item \textbf{Memory Efficiency:} By storing only the pooling indices (which require much less memory than full feature maps), SegNet significantly reduces the memory needed during training and inference.
    \item \textbf{Computational Simplicity:} Fewer parameters and a simpler decoder contribute to faster processing and more efficient training.
    \item \textbf{Boundary Preservation:} The method effectively preserves important spatial details, which is crucial for accurate tumor segmentation.
\end{itemize}

\paragraph{Limitations and Trade-Offs}
\label{sec:segnet_limitations}
While SegNet is efficient, it may produce slightly smoother segmentations compared to U-Net. The absence of skip connections means the decoder relies solely on pooling indices to reconstruct high-resolution details, which can sometimes lead to subtle inaccuracies in fine structures. Comparative studies \cite{kasar2024comparative,vedhasri2024segmentation} have shown that U-Net typically outperforms SegNet by achieving slightly higher Dice coefficients, though at a higher memory and computational cost.

\subsection{Preprocessing and Feature-Based Enhancements}
\label{sec:preprocessing_enhancements}

\subsubsection{Standard Preprocessing Techniques in MRI Segmentation}
\label{sec:standard_preprocessing}
Preprocessing is a vital step in preparing MRI data for segmentation. Common procedures include:
\begin{itemize}
    \item \textbf{Skull Stripping:} Removing non-brain tissues to focus on the region of interest.
    \item \textbf{Bias Field Correction:} Correcting intensity variations caused by inhomogeneities in the MRI scanner.
    \item \textbf{Intensity Normalization:} Standardizing the intensity values across images.
    \item \textbf{Data Augmentation:} Techniques such as rotation, scaling, and flipping to enhance the diversity of the training dataset.
\end{itemize}
These steps are crucial to reduce noise and ensure that segmentation models learn robust, clinically relevant features.

\subsubsection{Feature Detection and Augmentation}
\label{sec:feature_detection}
Recent advances have seen the integration of classical feature detectors (e.g., edge, texture, and gradient filters) into preprocessing pipelines. Some models have incorporated edge maps as additional input channels, guiding the network to focus on tumor boundaries. For example, studies such as that by Allah et al. \cite{allah2023edge} have demonstrated that supplementing CNN inputs with an edge-detected version of the MRI can result in improved segmentation performance. Such methods often compute edge maps using Sobel \cite{sobel1968isotropic} or Canny operators and then fuse these maps with the original intensity images to create composite (RGB-like) images.

\subsubsection{Impact of Preprocessing on Segmentation Quality}
\label{sec:preprocessing_impact}
A notable challenge in brain tumor segmentation is leveraging the complementary information present in different MRI modalities—commonly T1, T1 with contrast (T1c), T2, and FLAIR. Each modality provides distinct tissue contrasts: for example, FLAIR images emphasize edema while T1c highlights the active tumor core. An effective method to integrate these data is to fuse multiple modalities into a single three-channel (RGB) image, allowing standard CNN architectures pre-trained on color images to be utilized. Alqazzaz et al. \cite{alqazzaz2019brats} proposed training separate SegNet models for each modality, and then fusing the resulting feature maps (using element-wise maximum operations) to generate a composite representation that was subsequently classified by a decision tree. This strategy demonstrated improvements in segmentation performance by harnessing multi-modal information in a unified framework.

\begin{figure}[ht]
    \centering
    % Insert a diagram that illustrates an example of modality fusion
    \caption{An illustration of RGB-like image generation from multimodal MRI inputs, highlighting how each channel contributes complementary information for tumor segmentation.}
    \label{fig:modality_fusion}
\end{figure}

\subsubsection{Discussion on Preprocessing Impact and Feature Fusion}
\label{sec:preprocessing_discussion}
The literature consistently shows that while deep networks are capable of learning features directly from raw data, incorporating domain-specific information via preprocessing can further improve segmentation quality \cite{cheng2023feature}. In the case of brain tumor segmentation, enhanced edge maps or texture features can help differentiate tumor boundaries from surrounding tissue. Although adding these extra input channels may increase pre-processing complexity, the payoff in terms of segmentation accuracy is often significant. Several recent studies have confirmed that when preprocessing techniques emphasize contrast and edge details, segmentation accuracy—especially in difficult cases with low contrast or heterogeneous tissue—improves markedly.

\subsection{Detailed Analysis of SegNet-Based Studies}
\label{sec:segnet_studies}
This section focuses on recent research that is closely aligned with the thesis topic—studies that either implement SegNet for brain tumor segmentation or adopt similar feature-detection techniques.

\subsubsection{Modified SegNet with Hybrid Feature Extraction}
\label{sec:modified_segnet}
Palleti and Reddy \cite{palleti2025modified} developed a modified SegNet architecture that integrated texture-based feature extraction. In their approach, multi-modal MRI images (T1, T1c, T2, FLAIR) were first fused using an advanced image fusion technique, followed by denoising via a median filter. The resulting composite images were then input into a Modified SegNet that adopted a novel pooling strategy and further extracted texture features using Gabor filters \cite{gabor1946theory} and other descriptors. These features were fed into a classifier, which categorized tumor subregions with an overall classification accuracy of 98\% on the BraTS 2015 dataset. This study demonstrated that with proper preprocessing and feature fusion, the inherent limitations of standard SegNet could be overcome, rendering it highly effective for both segmentation and classification tasks.

\subsubsection{SegNet Optimized via Evolutionary Algorithms}
\label{sec:evolutionary_segnet}
Bidkar, Kumar, and Ghosh \cite{bidkar2022evolutionary} proposed an innovative approach by coupling SegNet with evolutionary optimization techniques. They applied Salp Swarm and Water Wave Optimization algorithms during the training phase to fine-tune SegNet's parameters, improving segmentation accuracy significantly. The optimized segmentation output was then processed by a Deep Belief Network (DBN) for tumor subregion classification on BraTS 2018 and 2020 datasets. The hybrid model achieved over 92\% accuracy and sensitivity, demonstrating that advanced training strategies can help SegNet rival other more complex architectures in performance, even though it is inherently more resource-efficient.

\subsubsection{Comparative Evaluation of U-Net and SegNet}
\label{sec:comparative_eval}
Several studies have compared the performance of U-Net and SegNet for brain tumor segmentation. Kasar et al. \cite{kasar2024comparative} conducted a direct comparison using a public brain MRI dataset. Their results indicated that while U-Net achieved a slightly higher average Dice coefficient (approximately 0.76) compared to SegNet (approximately 0.67), SegNet still showed promising results with notably lower computational and memory requirements. Vedhasri and V. \cite{vedhasri2024segmentation} similarly found that U-Net's skip connections led to more precise segmentation outcomes; however, the simplicity and efficiency of SegNet provided practical advantages for systems with hardware limitations. These comparative studies have informed the trade-offs considered in this thesis, highlighting that while maximum segmentation accuracy is valuable, considerations such as processing speed and resource consumption are also critical in clinical environments.

\subsubsection{3D-SegNet: Extension to Volumetric Data}
\label{sec:3d_segnet}
Zhou et al. \cite{zhou20193dsegnet} expanded the application of SegNet to 3D segmentation by replacing 2D operations with 3D convolutions and pooling. Their 3D-SegNet was applied to volumetric brain MRI data to capture spatial continuity across slices. The model demonstrated an ability to maintain high segmentation accuracy while operating more efficiently than traditional 3D U-Net frameworks. Additionally, multi-task learning approaches that combined segmentation with survival prediction or tumor grading have emerged. Although these methods primarily use networks similar to U-Net, some researchers have experimented with hybrid models incorporating SegNet in one pipeline segment, further emphasizing the potential utility of an efficient, feature-based SegNet in comprehensive diagnostic systems.

\subsubsection{GAN-SegNet and Adversarial Training Approaches}
\label{sec:gan_segnet}
Cui et al. \cite{cui2021gansegnet} integrated SegNet into a Generative Adversarial Network (GAN) framework, coining the term GAN-SegNet. In their model, a generator network—constructed in the style of SegNet—produced segmentation maps, and a discriminator network evaluated the realism of these maps. This adversarial setup pushed the generator toward producing more refined segmentations, achieving a Dice score greater than 0.90 on whole tumor segmentation tasks with BraTS 2018 data. Although GAN-SegNet was more complex than standard SegNet, the incorporation of adversarial training provided evidence that even resource-efficient architectures like SegNet could be enhanced through additional training paradigms.

\subsubsection{Hybrid and Ensemble Approaches}
\label{sec:ensemble_segnet}
Collectively, the studies reviewed illustrate a clear trend: while U-Net and its variants frequently achieve the highest segmentation accuracy, SegNet offers compelling benefits in terms of efficiency and flexibility. Early implementations of SegNet produced smoother output maps, but subsequent enhancements—ranging from modified loss functions to feature-enhanced preprocessing—have improved its performance considerably. The research suggests that when augmented with robust feature-detection and multi-modal fusion techniques, SegNet-based systems can serve as effective and resource-efficient solutions for brain tumor MRI segmentation.

\subsection{Synthesis and Rationale for a SegNet-Based, Feature-Detection Approach}
\label{sec:synthesis_rationale}
The literature indicates that despite the widespread adoption of U-Net for biomedical image segmentation, there is a substantial need for efficient models that perform well on resource-limited systems. SegNet, with its innovative use of pooling indices for upsampling, has demonstrated lower memory consumption and faster inference times. However, its basic version may fall short of U-Net in terms of boundary detail. Studies have shown that performance differences can be mitigated by integrating feature detection techniques—including edge, texture, and modality fusion—that enhance the input to the network.

The proposed approach—RGB Images Generation Based on Feature Detection Techniques: The Case of SegNet—builds upon these insights. By preprocessing MRI data to generate enhanced "RGB" images where additional channels carry explicit edge and texture information, it is possible to compensate for SegNet's lack of skip connections. This strategy not only boosts segmentation quality but also maintains the benefits of SegNet's efficiency. Moreover, given that many clinical settings contend with limited computing resources, the compact nature of SegNet makes it an attractive backbone when deployed in real-time or near-real-time diagnostic systems.

Ultimately, this integrated methodology—combining SegNet's efficient architecture with enhanced feature detection and multi-modal fusion—aligns with the evolving trends in brain tumor segmentation research. It strikes a balance between the high accuracy typically achieved by more complex models and the necessity for a lightweight, resource-efficient solution that can be feasibly used in clinical environments.

\subsection{Future Research Directions and Emerging Trends}
\label{sec:future_trends}
While significant advances have been made in brain tumor segmentation using deep learning, several challenges remain unresolved:
\begin{itemize}
    \item \textbf{3D Volumetric Segmentation:} Extending 2D segmentation approaches to fully 3D models continues to be challenging due to increased computational demand and memory constraints. Future research may explore hybrid 2D-3D models or efficient 3D implementations that maintain high accuracy with reasonable resource use.
    
    \item \textbf{Integration of Attention Mechanisms:} Attention-based modules have shown promise in guiding networks to focus on relevant regions. Future architectures may integrate attention within SegNet to recover some of the spatial detail lost by the absence of skip connections.
    
    \item \textbf{Unified Multi-Modal Fusion:} While some studies have successfully merged multiple MRI modalities into composite images, the optimal method for fusing modalities remains an open question. Research in this area could lead to more robust fusion techniques that further enhance segmentation performance.
    
    \item \textbf{Domain Adaptation and Generalization:} Variability among MRI scanners and protocols can affect model performance. Future work should address domain adaptation to ensure models generalize well across diverse clinical settings.
    
    \item \textbf{Adversarial and Hybrid Training Approaches:} The incorporation of GANs and other adversarial training techniques has shown promise in refining segmentation outputs. Such hybrid approaches warrant further exploration, particularly in balancing accuracy with computational cost.
\end{itemize}

\subsection{Conclusion}
\label{sec:lit_conclusion}
This literature review has traced the evolution of brain tumor segmentation methods from early, classical techniques to advanced deep learning approaches. The transition from manual and handcrafted methods to fully automated, data-driven segmentation has led to substantial improvements in efficiency and accuracy. While U-Net has often emerged as the state-of-the-art model due to its effective use of skip connections, SegNet represents a promising alternative because of its memory efficiency and streamlined architecture via pooling index transfer.

Studies that integrated feature-detection techniques—such as edge detection, texture filtering, and multi-modal fusion—demonstrated that augmenting input data with explicit, domain-specific cues can significantly enhance segmentation quality. These findings provide a strong rationale for the current thesis approach, which employs SegNet as a backbone augmented with feature-based preprocessing to generate RGB segmentation maps. Such an approach aims to deliver a lightweight yet accurate system for brain tumor segmentation, capable of operating under real-world clinical constraints.

As research continues to advance, combining efficient deep learning architectures with targeted preprocessing strategies offers great promise in achieving reliable, real-time segmentation of brain tumors. This review lays the foundation for further exploration by establishing that a SegNet-based, feature-detection method is well-suited to address both the technical and practical challenges inherent in brain tumor MRI segmentation.

\section{Key Background Terms}
\label{sec:key_terms}

Below is a list of key deep learning concepts and methods, explained simply and tailored to the context of segmenting brain tumors in MRI scans. We start with general ideas (like how CNNs work) and then zoom in to the specific techniques used in the SegNet model.

\subsection{Core Deep Learning Concepts}
\label{sec:core_dl_concepts}

\subsubsection{Convolutional Neural Networks (CNNs)}
\label{sec:cnn_definition}

A Convolutional Neural Network (CNN) is a type of deep learning model especially good at image analysis. Instead of looking at an entire image pixel-by-pixel, a CNN scans small patches with learnable filters to detect patterns. This process is similar to sliding a small window over a picture to find particular shapes or edges. CNNs were originally inspired by how the human visual cortex works, identifying simple features (like edges or textures) in early layers and more complex features (like shapes or objects) in later layers \cite{datacamp2023cnn}. 

In practical terms, a CNN learns to extract features automatically – for example, it might learn to detect the outline of a tumor or differences in texture between tumor tissue and healthy brain tissue, without needing a human to specify those features. CNNs have proven very effective in tasks like image classification, object detection, and importantly image segmentation (dividing an image into meaningful parts) \cite{datacamp2023cnn}. In the case of brain tumor segmentation, the CNN takes MRI images as input and learns to recognize patterns that distinguish tumor regions from normal brain, essentially learning what a tumor "looks like" across the various MRI modalities.

\textbf{Analogy:} Imagine a CNN as a team of detectors scanning an MRI slice. One detector might be looking for bright circular shapes (which could indicate tumor edema on a T2 image), another might look for irregular edges (tumor boundaries), and so on. Each detector filter slides over the image and lights up when it finds its particular pattern – much like highlighting all areas of an image that contain a certain texture. The network layers combine all these highlighted patterns to decide which pixels belong to tumor tissue. This approach helps the model autonomously learn visual features relevant to tumors, rather than relying on manual measurement of intensity or shape. It's a powerful advantage: the CNN learns the best features for the task by itself, which is why CNNs outperform many traditional methods in medical image analysis.

\subsubsection{Convolutional Filters and Feature Detection}
\label{sec:conv_filters}

A CNN's core strength lies in its convolutional filters – these are small matrices (e.g. 3x3 or 5x5) that slide over the image pixels. Each filter is like a tiny detector tuned to a specific feature. For instance, one filter may respond strongly to edges, another to a certain texture. In the first layers of a CNN, filters typically detect simple features such as horizontal or vertical edges, lines, or basic shapes \cite{medium2022cnn}. Deeper into the network, the filters' perceptions combine to detect complex features – for example, a pattern that might correspond to the rough shape of a tumor or the difference in texture between tumor core and surrounding edema. You can think of this as a hierarchy: early layers ask "Is there an edge here?" while later layers ask "Do these combined edges and textures look like part of a tumor?" \cite{datacamp2023cnn}.

This feature detection is crucial for brain tumor segmentation. Different brain tumor tissues have distinct appearances in MRI (for example, a tumor might have a bright rim in a contrast-enhanced T1 image or a halo of edema in FLAIR). The CNN's filters will learn to pick up on these cues. For example, a filter could learn the ring-like enhancement pattern often seen around certain tumor types, while another might learn the texture of normal brain tissue. By learning a rich set of features, the CNN can accurately distinguish tumor pixels from healthy pixels. In summary, convolutional filters enable the model to automatically detect relevant features in MRI scans – much like having a set of eyes, each looking for a specific visual clue of the tumor. This automated feature learning is why CNN-based methods are so powerful for medical image tasks.

\subsubsection{Pooling Layers (Downsampling)}
\label{sec:pooling_layers}

Pooling layers are components in CNNs that reduce the spatial size of feature maps, essentially compressing information while preserving the most important content. After convolutional layers have detected features, pooling summarises those features to make the representation more manageable and invariant to small shifts. The most common type is max pooling, which looks at a small region (e.g. 2x2 pixels) and retains only the maximum value in that region. This is like taking the strongest response in an area – essentially saying "if any part of this small patch had a feature, we'll assume the patch as a whole has that feature." The goal of pooling is to "pull out" the most significant signal from each region \cite{datacamp2023cnn}. As a result, the image gets down-sampled (shrunk), reducing detail but keeping the high-level patterns.

\textbf{Analogy:} Pooling is like zooming out on an image to get a simpler view \cite{medium2022cnn}. Imagine you have a very detailed map and you squint or step back; you lose some tiny details (like small streets) but you still see the major highways and cities. Similarly, after pooling, the network has a smaller, more abstract picture of the image – it's lost exact pixel details but retained the important features ("highways") that indicate structures in the image.

In brain tumor segmentation, pooling helps the network recognize the tumor in a way that's somewhat tolerant to slight differences in position or shape. For example, whether the tumor is slightly to the left or right, or whether the MRI is a bit shifted, the max pooling operation ensures that as long as a tumor feature is in a general area, it will be noted. However, pooling also means a loss of precise spatial detail – which is a trade-off. Too much pooling can make it hard to exactly outline the tumor's boundaries (since we "zoomed out" and lost fine detail). Segmentation networks like SegNet address this by having a decoder that later upsamples (brings the image back to full size) and recovers the spatial details using special techniques (more on that soon). In summary, pooling makes the CNN robust to small positional variations and reduces computation, but one must recover the detail later for pixel-perfect tumor maps. SegNet specifically remembers where the maximum came from (the pooling indices) so it can put features back in the right place during upsampling – a clever solution to pooling's loss of detail.

\subsubsection{Semantic Segmentation}
\label{sec:semantic_segmentation}

Semantic segmentation is the task of labeling every pixel in an image with a class (category) label. In other words, it doesn't just ask "What is in this image?" but rather "What is this pixel, and what about this pixel, and so on for all pixels?" \cite{ieee2023dnn}. The end result is typically a segmentation map or mask that has the same dimensions as the input image, where each pixel is color-coded or marked as belonging to a certain class. In our case, the classes could be "tumor" vs "non-tumor" (and possibly sub-classes like edema, necrosis, etc., if doing a more detailed tumor segmentation). Semantic segmentation is like painting the image with different colors for different tissues: e.g., all tumor pixels might be marked red and everything else transparent or another color. This is different from just detecting the tumor with a bounding box or a single label; segmentation gives a detailed outline of the tumor.

In brain tumor MRI analysis, semantic segmentation is exactly what we need – we want to pinpoint which pixels belong to the tumor. For instance, given an MRI slice of the brain, a segmentation model will output an image of the same size where each pixel has a label: 1 for tumor, 0 for background (for a simple two-class case). This fine-grained labeling is crucial for doctors to understand the tumor's size, shape, and location. It has applications in treatment planning (like surgery or radiotherapy), where knowing the precise boundary of a tumor can guide surgical margins or radiation contours. Semantic segmentation with CNNs replaces the labor-intensive process of manual contouring. Because every pixel is classified, metrics like tumor volume can be computed directly. However, it's also a challenging task – the model must be accurate at a very granular level. That's why specialized architectures (like encoder-decoder CNNs) are used, as they are designed to output dense pixel-wise predictions. Brain tumor segmentation networks take advantage of multiple MRI modalities and the powerful feature learning of CNNs to make this pixel-level classification accurate and robust.

\subsubsection{Encoder–Decoder Architecture}
\label{sec:encoder_decoder_arch}

An encoder-decoder architecture is a neural network design pattern commonly used for segmentation problems. It consists of two parts: an encoder that gradually downsamples the input (like a compressor), and a decoder that upsamples it back to the original resolution (like a reconstructor). You can imagine it as an hourglass shape or a funnel followed by an expanding cone. The encoder is typically a series of convolutional and pooling layers that take the large input image and squeeze it into a smaller, high-level representation. As it goes deeper, the encoder condenses the image information into important features while discarding extraneous details. By the end of the encoder, the network has a very condensed "essence" of the image (sometimes called a latent representation or feature map), which strongly indicates where important structures are, but in a low resolution. The decoder then takes this compact representation and builds the image back up, pixel by pixel, to the original size. It does this through upsampling operations (the inverse of pooling) and additional convolutions that refine the details. Essentially, the decoder learns how to take the encoder's abstract features (like "there is a round tumor-ish feature in this region") and translate them into a precise pixel map ("color these exact pixels as tumor").

In many segmentation CNNs (such as U-Net, SegNet, etc.), this encoder-decoder scheme is used because it combines the best of both worlds: the encoder captures context and what the image is globally (e.g. "there's a tumor somewhere in the top left"), and the decoder recovers spatial detail so that we know exactly where that tumor is at the pixel level \cite{guide2024segmentation}. In brain tumor segmentation, using an encoder-decoder network means the model can both understand the overall structure of the brain (and tumor location in context) and also output a full-size mask marking the tumor boundaries. The encoder might, for example, encode features that differentiate tumor tissue from normal tissue, while the decoder ensures that this information is output in the correct anatomical locations matching the MRI. Without a decoder, a CNN might only tell you presence/absence or roughly where something is (like classification or coarse heatmaps). The decoder makes sure you get a detailed segmentation map.

SegNet is one such encoder-decoder architecture for segmentation \cite{paperswithcode2023segnet}. Its encoder is based on a proven classification network (VGG16) that progressively reduces image size, and its decoder mirrors this process in reverse to produce a segmentation. This approach is particularly well-suited to medical imaging tasks like ours because we often need the final output to be the same size as the input (so that each output pixel aligns with a location in the MRI).

\subsection{SegNet Architecture for Semantic Segmentation}
\label{sec:segnet_architecture}

\subsubsection{Overview}
\label{sec:segnet_overview_section}

SegNet is a specific deep learning model that follows the encoder-decoder approach, designed explicitly for semantic segmentation tasks \cite{paperswithcode2023segnet}. It was originally introduced for segmenting scenes (like roads, buildings, etc., in driving images) but its design makes it very applicable to medical images as well.

\subsubsection{Encoder}
\label{sec:segnet_encoder}

The encoder of SegNet is essentially a series of convolutional layers and max-pooling layers (just like a standard CNN for image recognition). In fact, the original SegNet encoder is topologically identical to the convolutional part of the VGG16 network \cite{paperswithcode2023segnet} – a well-known image classification network. This encoder takes the input image (which could be a multi-channel MRI image, combining T1, T2, FLAIR, etc.) and processes it through layers of filters and pooling. As it goes deeper, the feature maps become smaller in spatial size but richer in information (encodes "what" is in the image, not "where" in precise detail). After the encoder stage, the image has been reduced to much lower resolution feature maps that strongly indicate the presence of various features (like "tumor texture detected here" or "edge of brain detected there").

\subsubsection{Decoder}
\label{sec:segnet_decoder}

The decoder is the novel part of SegNet. Its job is to take those low-resolution encoder feature maps and upsample them back to the original image size, producing a dense segmentation output. But simply upsampling (like interpolation) could blur things; SegNet does something clever to improve this: it uses the pooling indices from the encoder. During each encoder pooling (downsampling) step, SegNet remembers the indices (locations) of the maximum values within each pooling window \cite{paperswithcode2023segnet}. The decoder then takes a correspondingly sized upsampling window and places the feature values back into their original positions using those indices (this is often called unpooling with indices). This effectively transfers the rough location information that was lost during pooling back into the decoder. After this "guided" upsampling, each decoder stage applies convolutional filters to refine the features (since unpooling with indices creates sparse feature maps) \cite{fezan2023segmentation}. By the final decoder layer, we get feature maps at full image resolution. These are then fed into a pixel-wise classification layer (typically a softmax) that assigns a class label to each pixel \cite{fezan2023segmentation}.

\subsubsection{Intuition Behind Unpooling}
\label{sec:unpooling_intuition}

In simpler terms, SegNet's decoder can be imagined as taking the encoder's compressed knowledge ("there is something of interest in these general areas") and using the stored indices as pointers to paint that knowledge back onto a high-resolution canvas. Each decoder layer adds more detail, guided by where the encoder found strong activations. The end result is an output the same size as the input MRI, where each pixel now has a probability or label of being tumor or not. For example, if the encoder discovered a region with high tumor-like features, the decoder will upsample that and mark the corresponding pixels in the output mask as tumor.

\subsubsection{Pooling-Index Transfer}
\label{sec:pooling_index_transfer}

The key innovation in SegNet is this pooling index transfer mechanism \cite{paperswithcode2023segnet}. It makes the network more memory-efficient and effective: instead of storing entire feature maps from the encoder (as some other methods do with skip connections), SegNet only stores the index of the strongest feature in each region. This is much lighter but still gives the decoder what it needs to reconstruct sharp boundaries. In tasks like brain tumor segmentation, preserving boundary information is critical – we need precise outlines. SegNet's approach of remembering where the strongest edge/feature was (via indices) helps the decoder draw more exact edges of the tumor. The decoder essentially says "I know one important feature was exactly at this pixel (because of the saved index), so when I upsample, I'll put a feature there." This yields sharper segmentation outputs than a naive upsampling would.

\subsubsection{Summary}
\label{sec:segnet_summary}

To sum up, SegNet is an encoder-decoder CNN specialized for segmentation, which stores pooling indices in the encoder and uses them in the decoder to get accurate upsampling \cite{paperswithcode2023segnet}. It has no fully connected layers, which means it's entirely convolutional and can handle images of variable size and use the full context of the image. This architecture is well-suited to segmentation problems where you need pixel-level accuracy but also want to benefit from the abstraction that pooling provides.

\subsection{Additional Key Concepts}
\label{sec:additional_concepts}

\subsubsection{Pooling Index Transfer (Detailed)}
\label{sec:pooling_index_detailed}

Pooling index transfer refers to the method SegNet uses to perform upsampling in the decoder by reusing the information from the encoder's pooling layers. Let's break that down: in each encoder layer, when we do max pooling, we not only get a pooled output, but we can also record where the maximum came from in each pooling window. Think of it like leaving breadcrumbs. Each breadcrumb (index) tells us "out of this 2x2 block, the top-right was the brightest spot," for example. SegNet collects these breadcrumbs (indices) for every pooling layer in the encoder.

When it's time to decode (upsample), instead of just blindly stretching the low-resolution feature map, SegNet places each feature back to its original pooled location according to the saved index \cite{paperswithcode2023segnet}. This operation is sometimes called unpooling. The result of unpooling with indices is a sparse feature map (mostly zeros, with feature values at certain locations). This is then followed by a convolution to fill in and smooth the gaps, effectively "joining the dots" and producing a dense map \cite{fezan2023segmentation}. By doing this, the decoder has effectively reused the encoder's knowledge of where important features were. It's as if the encoder said, "I found a strong edge at these coordinates," and the decoder, later on, uses that exact coordinate to help reconstruct the edge in the output.

Why is this helpful? Because one big challenge in segmentation is maintaining precise boundary details. Pooling by nature loses some spatial precision – we know a feature existed in a 2x2 region, but not exactly where. The pooling index transfer gives that precision back. It has a memory of the exact pixel position of maximal response in each small region \cite{fezan2023segmentation}. Therefore, boundaries and shapes can be more sharply reconstructed. In brain tumor segmentation, this is very important: the difference between including or excluding a single pixel at the edge could mean a lot when estimating tumor size. SegNet's indices help ensure that if the encoder spotted the tumor's boundary, the decoder can mark that same boundary location.

Another advantage is efficiency. Storing indices (basically one number per pooling window) is cheap compared to storing entire feature maps (which is what U-Net does via skip connections). This makes SegNet use less memory. It was one of the motivations for the SegNet design – to be memory efficient yet accurate \cite{fezan2023segmentation}. This efficiency can be important when working with large 3D medical images or limited GPU memory.

\textbf{Analogy:} Suppose you are reading a condensed version of a textbook where many details have been left out, but it says "refer to page X for details on Y." If you have the original textbook, you can go to page X to get the detailed info. In SegNet, the encoder outputs are like the condensed version – they have pointers (indices) to where the detail was. The decoder is like having the reference to go back and fill in the details at the right spots. The result is you don't lose the critical information. In practice, pooling index transfer in SegNet allows it to produce segmentation maps that are both precise and obtained efficiently, which is why even though it doesn't explicitly pass along all feature maps (like some other networks), it still performs well. For our MRI tumor task, it means the network can delineate the tumor region with pixel-level accuracy, guided by the indices captured from the high-resolution encoder features.

\subsubsection{MRI Modalities (T1, T2, FLAIR)}
\label{sec:mri_modalities}

In brain tumor segmentation, we often leverage multiple types of MRI scans, known as modalities. The common ones are T1, T2, and FLAIR (Fluid-Attenuated Inversion Recovery). Each of these MRI modalities is like looking at the brain with a different filter or lens – they emphasize different biological aspects, which helps in distinguishing tumor tissue from normal tissue. Here's a simple rundown of each \cite{nature2023mri}:

\paragraph{T1-weighted MRI}
Think of T1 as giving a view of the brain's basic anatomy. On T1 images, fat appears bright and fluid (like cerebrospinal fluid, CSF) appears dark. T1 provides clear detail of structures – for example, gray matter vs white matter in the brain can be distinguished. Tumors on plain T1 might not always be obvious, but when a contrast agent (gadolinium) is given, tumor tissue often lights up (enhances) on T1. In our context, if we mention T1 without contrast, it shows the structure; with contrast (sometimes denoted T1c or T1-Gd), the tumor's blood-brain barrier breakdown areas will glow. In general, T1 helps observe normal structures and, with contrast, highlights abnormal lesions. It gives a kind of "baseline" view of the brain \cite{researchgate2023modalities}.

\paragraph{T2-weighted MRI}
T2 images make fluid appear bright. This means CSF is white on T2. Importantly, many tumors contain edema (swelling with fluid) around them or inside them, which will also appear bright on T2. T2 is excellent for seeing anything with high water content – swelling, cystic parts of tumors, etc. Clinicians often say that T2 is good for finding lesions, because an abnormality (like a tumor or inflammation) often stands out as a bright spot against darker normal brain tissue \cite{researchgate2023modalities}. In brain tumor segmentation, a T2 image will usually show the edema around the tumor as a bright halo. So T2 helps locate the general region of a tumor and associated swelling, albeit with less anatomical detail than T1.

\paragraph{FLAIR (Fluid-Attenuated Inversion Recovery)}
FLAIR is a special kind of T2-weighted image where free fluid (like pure fluid in ventricles or CSF spaces) is made dark. In other words, it attenuates (removes) the signal from fluid. Why do that? Because it helps us see lesions that are next to fluid-rich areas. For example, a tumor near the ventricles (which are filled with CSF) might be hard to see on regular T2 since both the tumor edema and the CSF are bright. On FLAIR, the CSF fluid is suppressed (dark), but the edema in the tissue is still bright, so the lesion stands out more clearly \cite{researchgate2023modalities}. FLAIR is great for seeing abnormal hyperintensities in the brain parenchyma (brain tissue) without the distraction of bright CSF. In tumor cases, FLAIR is often the best at delineating the extent of edema/infiltrative tumor, because it shows the full abnormal signal in the tissue (everything that's not pure fluid).

Each modality provides a unique piece of the puzzle \cite{nature2023mri}. In segmentation, especially with modern challenges like BraTS (Brain Tumor Segmentation Challenge), they provide all modalities (T1, T1c, T2, FLAIR) to the algorithm, because combining them gives the full picture of the tumor. In our scenario, using T1, T2, and FLAIR means our CNN (like SegNet) can learn from all three "views" of the brain. For instance, it might use FLAIR to catch the full extent of the tumor and edema, T2 to confirm the presence of edema and fluid content, and T1 (with contrast) to pinpoint the active tumor core. When we feed these into the network, we typically stack them as separate channels (just like a color image has R, G, B channels, here we have a T1 channel, a T2 channel, etc.). The network's filters then learn from combinations of these modalities.

\textbf{Relatable idea:} If you were examining a suspicious spot in the brain, a T1 might be like a normal light photograph, T2 like an infrared view highlighting water, and FLAIR like a special high-contrast view that suppresses glare (fluid). A doctor uses all these sequences together to make a diagnosis; similarly, a CNN benefits from the complementary information. MRI modalities are thus essential inputs for brain tumor segmentation models. They increase the network's ability to differentiate tumor tissue (which might look different in each modality) from normal tissue. Our SegNet model will exploit the differences – for example, learning that "if something is bright on T2 and FLAIR but not on T1, it's likely edema around a tumor."

\subsubsection{Dice Score}
\label{sec:dice_score}

The Dice score (also called Dice Similarity Coefficient) is a metric to evaluate segmentation accuracy by comparing the overlap between the predicted segmentation and the ground truth (the manual, correct segmentation). It's a number between 0 and 1, where 1 indicates perfect overlap (the segmentation exactly matches the truth) and 0 indicates no overlap at all \cite{oecd2023dice}. In formula terms, Dice = 2 * (Area of Overlap) / (Area of Prediction + Area of Ground Truth). You can think of it as a measure of how well two shapes match. If the CNN marks almost the same region for the tumor as the radiologist did, the Dice score will be very high (close to 1). If it misses a lot or marks extra areas, the Dice score will be lower.

For example, suppose in an MRI slice the true tumor region has 100 pixels. If our SegNet correctly identified 90 of those and also mistakenly added 10 pixels that aren't tumor, the Dice score would be 2 * 90 / (100 + 100) = 0.90. If it somehow exactly matched all 100 and added none extra, it'd be 1.0 (perfect). But if it only found 50 of them and added 50 wrong ones (so overlap 50, our prediction 100, truth 100), Dice = 2*50/(100+100) = 0.5. Thus, Dice balances sensitivity and precision in one number.

In medical imaging, Dice score is very popular for segmentation evaluation because it directly measures the agreement between two areas. A high Dice means the algorithm's tumor mask is almost exactly what the doctor drew – which is what we want. It's especially useful when the size of the structure varies, or in imbalanced cases (like tumor vs background, where background is much larger; Dice focuses on the region of interest overlap). When we report results for brain tumor segmentation, we often say something like "our model achieved a Dice score of 0.85" meaning on average an 85\% overlap with expert-defined tumor regions – a strong result. We aim for high Dice scores as a sign that SegNet is accurately capturing the tumor. (Sometimes other metrics like Jaccard index or voxel-wise accuracy are used, but Dice is the go-to in medical segmentation).

In preparation for an oral defense or presentation, you can explain Dice in a non-technical way: "It's like how much the computer's outlined tumor matches the doctor's outlined tumor. 100\% means they completely agree; lower percentages mean there's disagreement or error." It provides an intuitive gauge of segmentation quality. Typically, one might show a few examples of segmented images alongside saying "and quantitatively, we measured a Dice score of X to assess accuracy." Knowing that Dice is overlap-based will help answer any questions on why we chose it (it's standard and meaningful for overlap tasks) \cite{oecd2023dice}.

\subsection{Why SegNet?}
\label{sec:why_segnet}

Finally, why did we choose SegNet as our model for this brain tumor segmentation task? There are several compelling reasons that make SegNet a good fit:

\begin{itemize}
    \item \textbf{Pixel-Wise Segmentation:} SegNet was designed specifically for semantic segmentation, so it outputs a label for \emph{every} pixel rather than just a bounding box. That aligns perfectly with our goal of classifying each MRI pixel as tumor vs. non-tumor.
    
    \item \textbf{Encoder–Decoder Efficiency:}
    \begin{itemize}
        \item \emph{Encoder:} Uses a pre-trained VGG16 convolutional backbone, giving strong feature extraction out-of-the-box.
        \item \emph{Decoder:} Lightweight un-pooling based on stored pooling indices instead of heavy skip connections, keeping memory and compute footprints low.
    \end{itemize}
    This efficiency lets us train on full-resolution MRI slices (or even 3D patches) and on multiple modalities without running out of GPU memory.
    
    \item \textbf{Boundary Preservation via Pooling Indices:} By saving only the locations of max-pool activations in the encoder and "unpooling" with those indices in the decoder, SegNet reconstructs sharp edges (critical for irregular tumor boundaries) without a naive interpolation step.
    
    \item \textbf{Multi-Modal Flexibility:} Being fully convolutional, SegNet accepts images with any number of channels. We simply extend its first conv-layer to 3 (or 4) channels for T1, T2, FLAIR (±T1c), letting the network learn cross-modal features in a single pass.
    
    \item \textbf{Proven in Practice:} SegNet has been successfully applied to multiple segmentation domains—including medical imaging—and comes with abundant reference implementations and training tips. This community support accelerates our development and debugging.
    
    \item \textbf{Simplicity: No Dense Layers:} SegNet omits fully connected layers entirely, keeping parameter counts lower and reducing overfitting risk on our relatively small MRI dataset. Its fully-convolutional nature also means flexible input sizes and seamless patch-to-full-image inference.
\end{itemize}

In summary, SegNet strikes a balance between accuracy (thanks to its pooling-index unpooling) and efficiency (few parameters, low memory), making it an ideal backbone for multi-modal, pixel-level brain tumor segmentation.

We chose SegNet for brain tumor segmentation because it strikes a good balance between accuracy and efficiency, and it aligns well with the needs of the problem. It provides a clear way to get pixel-level predictions and is known to perform well in segmentation benchmarks. Its pooling index mechanism is particularly useful for retaining the fine details of tumor boundaries \cite{fezan2023segmentation}, which is a top priority in medical imaging. By using SegNet, we leveraged a well-established architecture, adapted it to our multi-modal MRI input, and benefited from its strengths to achieve reliable tumor segmentation results.

Ultimately, in our presentation or defense, we can confidently explain that SegNet was an appropriate choice because it was built for exactly this kind of task. It allowed us to incorporate domain knowledge (multiple MRI sequences) and produce outputs that can be directly evaluated against expert annotations (using metrics like Dice score). The architecture's design choices (like encoder-decoder with index unpooling) can be justified in terms of the problem requirements (need for context + detail). And as we've shown through definitions and analogies above, even someone without a deep ML background can appreciate that SegNet takes an image, simplifies it to understand it, then reconstructs a detailed mask of the tumor – which is precisely what we want a brain tumor segmentation solution to do.

\textbf{Sources:} The explanations above were informed by key references in deep learning and medical imaging, including definitions of CNNs \cite{datacamp2023cnn,medium2022cnn}, the concept of semantic segmentation \cite{ieee2023dnn}, descriptions of the SegNet architecture and its pooling index approach \cite{paperswithcode2023segnet,fezan2023segmentation}, and insights on MRI modalities for brain tumors \cite{nature2023mri}, among others. Each concept has been presented in simplified terms to aid understanding while remaining rooted in these established sources.

\section{Implementation Progress}
\label{sec:implementation_progress}

\begin{itemize}
    \item \textbf{Model Definition:} \texttt{models/segnet.py} – custom PyTorch SegNet with 4-channel input, VGG16\_bn encoder, max-unpooling decoder.
    \item \textbf{Dataset Loader:} \texttt{src/dataset.py} – NiBabel-based BraTS slice indexing, per-slice normalization, label remapping.
    \item \textbf{Inference Script:} \texttt{src/inference.py} – one-sample demo, overlay visualization of FLAIR + predicted mask.
    \item \textbf{Training Loop:} \texttt{src/train.py} – combined CrossEntropy + Dice loss, Adam optimizer, checkpointing, ROCm diagnostics.
\end{itemize}

\noindent\textit{Current status: all components compile and run on CPU; pending GPU driver and CUDA setup before full end-to-end training can proceed.}